\section{\texttt{src/evaluation/metrics.py}}

Tutte le metriche del progetto, implementate da zero.

\begin{nota}
\textbf{Perché non usare \texttt{scipy.stats}.} Sono poche righe, e le convenzioni --- gestione
dei pari merito, ordine degli argomenti di KL, denominatore del false positive rate --- sono
esattamente ciò che va capito e saputo spiegare. Implementarle è parte del progetto, non una
perdita di tempo.
\end{nota}

\subsection{Le cinque metriche intrinseche}

\begin{center}
\begin{tabular}{l l l l}
\toprule
\rowcolor{tblHead}\textcolor{white}{\textbf{Metrica}} & \textcolor{white}{\textbf{Ground truth}} & \textcolor{white}{\textbf{Intervallo}} & \textcolor{white}{\textbf{Direzione}} \\
\midrule
CC & heatmap continua & $[-1, 1]$ & alto = meglio \\
\rowcolor{tblAlt} SIM & heatmap continua & $[0, 1]$ & alto = meglio \\
KL & heatmap continua & $[0, \infty)$ & \textbf{basso = meglio} \\
\rowcolor{tblAlt} NSS & fissazioni discrete & $(-\infty, \infty)$ & alto = meglio \\
AUC-Judd & fissazioni discrete & $[0, 1]$ & alto = meglio \\
\bottomrule
\end{tabular}
\end{center}

\begin{attenzione}
\textbf{KL va nella direzione opposta} rispetto a tutte le altre. Nelle tabelle di confronto va
segnalato esplicitamente, altrimenti si legge il risultato al contrario.
\end{attenzione}

\begin{nota}
\textbf{Perché cinque metriche e non una.} Ognuna è cieca a qualcosa. CC ignora dove sta il picco
purché l'andamento generale corrisponda; NSS guarda solo i punti fissati e ignora tutto il resto;
AUC-Judd è invariante a qualsiasi trasformazione monotona, quindi conta l'ordinamento ma non
l'intensità dei picchi.

\medskip
Un modello che produce una mappa perfetta ma spostata di dieci pixel avrà CC ancora alto e NSS
crollato. Riportare una sola metrica significa nascondere metà del comportamento.
\end{nota}

\subsection{Funzioni di supporto}

\begin{idebox}
\begin{pycode}
def _to_distribution(saliency_map):
    m = saliency_map.astype(np.float64)
    m = m - m.min()
    total = m.sum()
    if total < EPS:
        return np.full_like(m, 1.0 / m.size)
    return m / total
\end{pycode}
\end{idebox}

Riscala la mappa perché la somma dei valori sia 1, rendendola una distribuzione di probabilità.
Serve a SIM e KL.

Il cast a \pyin{float64} non è pedanteria: KL somma logaritmi su decine di migliaia di pixel, e
in \pyin{float32} l'errore di arrotondamento si accumula in modo visibile.

Il caso degenere (mappa costante) restituisce la distribuzione uniforme, l'unica sensata.

\begin{idebox}
\begin{pycode}
def _standardize(saliency_map):
    m = saliency_map.astype(np.float64)
    std = m.std()
    if std < EPS:
        return np.zeros_like(m)
    return (m - m.mean()) / std
\end{pycode}
\end{idebox}

Z-score: media 0, deviazione standard 1. Rende confrontabili mappe prodotte su scale diverse ed è
il cuore di CC e NSS.

\subsection{\texttt{cc} --- correlazione di Pearson}

\begin{idebox}
\begin{pycode}
def cc(saliency_map, ground_truth_map):
    _check_shape(saliency_map, ground_truth_map)
    return float((_standardize(saliency_map) * _standardize(ground_truth_map)).mean())
\end{pycode}
\end{idebox}

\begin{nota}
\textbf{Perché una riga basta.} La correlazione di Pearson è
$\text{cov}(S,G) / (\sigma_S \sigma_G)$. Standardizzando entrambe le mappe, le deviazioni standard
diventano 1 e le medie 0: la formula si riduce alla media del prodotto. Non serve implementare
covarianza e deviazioni separatamente.
\end{nota}

Conseguenza dell'invarianza: una mappa e la stessa mappa moltiplicata per 10 hanno CC $= 1$. Il
modello non deve essere calibrato, solo corretto nell'andamento.

\subsection{\texttt{sim} --- histogram intersection}

\begin{idebox}
\begin{pycode}
s = _to_distribution(saliency_map)
g = _to_distribution(ground_truth_map)
return float(np.minimum(s, g).sum())
\end{pycode}
\end{idebox}

Entrambe le mappe diventano distribuzioni a somma 1; si somma il minimo punto per punto: quanta
massa di probabilità è condivisa. Con distribuzioni identiche il minimo coincide con la mappa
stessa, e la somma vale 1.

\subsection{\texttt{kl\_divergence}}

\begin{idebox}
\begin{pycode}
support = g > 0
if not support.any():
    return 0.0
return float(np.sum(g[support] * np.log(g[support] / np.maximum(s[support], EPS))))
\end{pycode}
\end{idebox}

Tre dettagli, tutti emersi in fase di test:

\begin{description}[leftmargin=1.2em,style=nextline]
\item[La maschera \pyin{g > 0}] Dove la ground truth è nulla, il termine $g \log(g/s)$ vale 0 per
definizione (limite di $x\log x$ per $x \to 0$). Escludere quei punti è matematicamente corretto
e più rapido.

\item[\pyin{np.maximum(s, EPS)} e non \pyin{s + EPS}] Differenza sostanziale. Sommare
\pyin{EPS} altera il rapporto \emph{ovunque}, introducendo un errore di $-\varepsilon$ per pixel
che si accumula: su una mappa da 4800 pixel, KL di una distribuzione con se stessa dava
$-4.7 \times 10^{-9}$ invece di 0. Con \pyin{maximum}, \pyin{s} viene usato esattamente dove è
significativo, ed \pyin{EPS} interviene solo dove \pyin{s} è nullo.

\item[L'ordine degli argomenti] KL è \textbf{asimmetrica}: $\text{KL}(a\|b) \neq \text{KL}(b\|a)$.
La convenzione della letteratura è KL(ground truth $\|$ predizione), che punisce severamente il
modello che assegna saliency quasi nulla dove la ground truth ne ha molta --- nell'uso reale
mancare un oggetto importante è più grave che segnalarne uno di troppo.
\end{description}

\subsection{\texttt{nss}}

\begin{idebox}
\begin{pycode}
fixations = fixation_map > 0
if not fixations.any():
    return float("nan")
return float(_standardize(saliency_map)[fixations].mean())
\end{pycode}
\end{idebox}

Standardizza la mappa predetta e ne prende il valore medio \emph{esattamente sui punti fissati}.
L'indicizzazione booleana \pyin{[fixations]} estrae solo quei valori.

Il valore 0 corrisponde alla prestazione casuale: le fissazioni cadono su punti di saliency
media. Valori positivi indicano che cadono su zone più salienti della media.

Senza fissazioni si restituisce \pyin{nan} invece di 0: 0 significherebbe ``prestazione casuale''
e falserebbe le medie sul dataset.

\subsection{\texttt{auc\_judd}}

La più elaborata. Interpreta la saliency map come un classificatore binario: al variare della
soglia, quante fissazioni vengono catturate rispetto a quanta immagine viene selezionata?

\begin{idebox}
\begin{pycode}
fixation_values = np.sort(s[fixations])[::-1]
all_values = np.sort(s.ravel())
n_fixations = fixation_values.size
n_pixels = all_values.size
n_non_fixations = n_pixels - n_fixations

for i, threshold in enumerate(fixation_values):
    above = n_pixels - np.searchsorted(all_values, threshold, side="left")
    tp_rates.append((i + 1) / n_fixations)
    fp_rates.append((above - (i + 1)) / n_non_fixations)
\end{pycode}
\end{idebox}

\subsubsection{Le soglie}

Si usano i valori di saliency \emph{nei punti fissati}, ordinati in modo decrescente. Sono gli
unici punti in cui il true positive rate cambia: fra due valori consecutivi il numeratore resta
costante, quindi bastano a descrivere tutta la curva.

\subsubsection{La ricerca binaria}

\pyin{np.searchsorted} su un array già ordinato costa $O(\log n)$. L'alternativa ingenua
(\pyin{(s >= threshold).sum()}) costa $O(n)$ e va ripetuta per ogni fissazione: su un'immagine
800$\times$600 con 18 fissazioni la differenza è già misurabile.

\subsubsection{Il denominatore del false positive rate}

\begin{nota}
\pyin{(above - (i+1)) / n_non_fixations}, non \pyin{above / n_pixels}. Le fissazioni sopra soglia
sono \emph{veri} positivi: contarle anche come falsi positivi gonfierebbe l'asse orizzontale
della curva ROC. Si sottraggono dal numeratore e si escludono dal denominatore.
\end{nota}

\subsubsection{Un limite matematico da conoscere}

\begin{attenzione}
Con una \textbf{sola} fissazione la curva ROC ha soltanto tre punti --- $(0,0)$,
$(\text{fp},1)$, $(1,1)$ --- e la sua area vale $1 - \text{fp}/2$. Il minimo è quindi
\textbf{0.5}, raggiunto quando $\text{fp} = 1$ (fissazione sul valore minimo della mappa).

\medskip
Servono almeno due fissazioni perché l'AUC possa scendere sotto 0.5 e segnalare una predizione
peggiore del caso. Questo comportamento è emerso scrivendo i test: l'aspettativa iniziale
(``fissazione sul minimo $\to$ AUC vicina a 0'') era matematicamente impossibile. Il test è stato
corretto, non il codice --- ma solo dopo aver verificato su carta quale dei due sbagliasse.
\end{attenzione}

\subsection{Le metriche di ranking}

\subsubsection{\texttt{rank\_data} --- i ranghi medi}

\begin{idebox}
\begin{pycode}
order = np.argsort(values, kind="mergesort")
sorted_values = values[order]

i = 0
while i < n:
    j = i
    while j + 1 < n and sorted_values[j + 1] == sorted_values[i]:
        j += 1
    ranks[order[i:j + 1]] = (i + j) / 2.0 + 1.0
    i = j + 1
\end{pycode}
\end{idebox}

Converte valori in ranghi assegnando la \textbf{media} in caso di pari merito:
$[10, 20, 20, 40] \to [1, 2.5, 2.5, 4]$. I due valori uguali occuperebbero le posizioni 2 e 3:
entrambi ricevono 2.5.

Il ciclo interno individua l'estensione di ogni gruppo di valori identici; \pyin{order[i:j+1]}
riporta il rango medio alle posizioni originali.

\pyin{kind="mergesort"} garantisce un ordinamento \emph{stabile}: a parità di valore l'ordine
originale è preservato, rendendo il risultato deterministico.

\subsubsection{\texttt{spearman}}

\begin{idebox}
\begin{pycode}
rp = rank_data(predicted)
rt = rank_data(reference)
rp = rp - rp.mean()
rt = rt - rt.mean()

denominator = np.sqrt((rp ** 2).sum() * (rt ** 2).sum())
if denominator < EPS:
    return float("nan")
return float((rp * rt).sum() / denominator)
\end{pycode}
\end{idebox}

È la correlazione di Pearson calcolata sui \emph{ranghi} invece che sui valori. Cattura qualsiasi
relazione monotona: $[1,2,3,4]$ e $[10, 200, 3000, 40000]$ hanno $\rho = 1$ pur non essendo
linearmente correlati.

Si usa la formula generale e non la scorciatoia $1 - 6\sum d^2 / n(n^2-1)$, perché quest'ultima
vale solo in assenza di pari merito.

\subsubsection{\texttt{kendall\_tau} --- e un errore istruttivo}

\[
\tau_b = \frac{C - D}{\sqrt{(n_0 - n_1)(n_0 - n_2)}}
\]

dove $n_0$ è il numero totale di coppie, $n_1$ le coppie in pari merito nella predizione, $n_2$
quelle in pari merito nella ground truth.

\begin{idebox}
\begin{pycode}
for i in range(n - 1):
    for j in range(i + 1, n):
        dp = predicted[i] - predicted[j]
        dt = reference[i] - reference[j]

        if dp == 0:
            ties_predicted += 1
        if dt == 0:
            ties_reference += 1

        if dp != 0 and dt != 0:
            if dp * dt > 0:
                concordant += 1
            else:
                discordant += 1

n_pairs = n * (n - 1) // 2
denominator = np.sqrt((n_pairs - ties_predicted) * (n_pairs - ties_reference))
\end{pycode}
\end{idebox}

\begin{attenzione}
\textbf{L'errore scoperto dai test.} La prima versione usava
$\sqrt{(C + D + n_1)(C + D + n_2)}$, che sembra equivalente ma non lo è.

\medskip
$n_0 - n_1$ conta le coppie \emph{non} in pari merito nella predizione, mentre $C + D + n_1$
conta quelle decise più quelle pari. Le due espressioni coincidono quando nessuna coppia è in
pari merito in \textbf{entrambe} le sequenze --- e anche quando differiscono, il \emph{prodotto}
dei due fattori resta identico perché i termini si scambiano.

\medskip
Risultato: l'errore era invisibile in tutti i test casuali e anche nel confronto con
\texttt{scipy}. È emerso solo sul caso limite \pyin{kendall_tau([5,5,5], [5,5,5])}, dove la
formula corretta dà $\sqrt{0 \cdot 0} = 0$ e quindi \pyin{nan} (``metrica non definita''), mentre
quella sbagliata dava $0/3 = 0$ (``nessuna correlazione''). Significati completamente diversi.
\end{attenzione}

\subsubsection{\texttt{top\_k\_accuracy}}

\begin{idebox}
\begin{pycode}
true_best = int(np.argmax(reference))
top_k_predicted = np.argsort(predicted)[::-1][:k]
return float(true_best in top_k_predicted)
\end{pycode}
\end{idebox}

Metrica binaria per singola immagine: 1 se l'oggetto più importante secondo la ground truth è nei
primi $k$ posti della predizione. La media sul dataset diventa una percentuale.

\begin{nota}
\textbf{Perché Top-1 conta più delle altre in questo progetto.} Lo scenario applicativo è
l'assistenza visiva: un sistema che descrive la scena a chi non la vede bene. Sbagliare l'ordine
fra il terzo e il quarto oggetto è irrilevante; sbagliare \emph{quale sia} l'oggetto principale
compromette l'intera descrizione.
\end{nota}

\section{\texttt{src/ranking/objects.py}}

Il passaggio finale: dalle saliency map a una classifica di oggetti.

\subsection{L'architettura del confronto}

\begin{center}
\begin{tabular}{l l}
\toprule
\rowcolor{tblHead}\textcolor{white}{\textbf{Percorso}} & \textcolor{white}{\textbf{Produce}} \\
\midrule
saliency map + bounding box & ranking \textbf{predetto} \\
\rowcolor{tblAlt} click annotati + bounding box & ranking \textbf{di riferimento} \\
\bottomrule
\end{tabular}
\end{center}

I due percorsi condividono gli stessi box e differiscono solo per la fonte dell'importanza. Il
confronto avviene con Spearman, Kendall e Top-1.

\subsection{\texttt{crop\_box}}

\begin{idebox}
\begin{pycode}
x1 = int(np.clip(x1, 0, width))
x2 = int(np.clip(x2, 0, width))
y1 = int(np.clip(y1, 0, height))
y2 = int(np.clip(y2, 0, height))

if x2 <= x1 or y2 <= y1:
    return np.empty((0, 0), dtype=saliency_map.dtype)

return saliency_map[y1:y2, x1:x2]
\end{pycode}
\end{idebox}

Il \pyin{np.clip} gestisce i box che sforano i bordi: i detector li producono regolarmente quando
un oggetto è parzialmente fuori inquadratura. Ricondurli dentro è più utile che sollevare
un'eccezione a metà valutazione.

Il controllo successivo intercetta i box degeneri e quelli interamente fuori: dopo il clip
diventerebbero intervalli vuoti, e \pyin{.mean()} su un array vuoto emette un warning
restituendo \pyin{nan}.

\subsection{\texttt{score\_box} --- le tre strategie}

\begin{center}
\begin{tabular}{l p{5.4cm} p{4.4cm}}
\toprule
\rowcolor{tblHead}\textcolor{white}{\textbf{Strategia}} & \textcolor{white}{\textbf{Domanda a cui risponde}} & \textcolor{white}{\textbf{Dipendenza dall'area}} \\
\midrule
\pyin{mean} & Quanto è saliente in media? & nessuna \\
\rowcolor{tblAlt} \pyin{max} & Contiene il punto più saliente? & nessuna, ma fragile \\
\pyin{sum} & Quanta attenzione totale raccoglie? & \textbf{forte} \\
\bottomrule
\end{tabular}
\end{center}

\begin{nota}
\textbf{Perché implementarle tutte e tre.} Non sono alternative fra cui scegliere a priori: il
loro confronto è una delle analisi previste dal progetto. \pyin{sum} è inclusa proprio per
rendere \emph{misurabile} il bias dell'area, non perché sia sbagliata in assoluto --- se
l'obiettivo fosse ``quale oggetto occupa più attenzione complessiva'', sarebbe la scelta corretta.

\medskip
Il risultato non è banale: un modello ottimo su tutte le metriche intrinseche può fallire il
ranking finale se si sceglie la strategia sbagliata. La qualità della saliency map non basta:
conta anche come la si traduce in punteggi per oggetto.
\end{nota}

Il default è \pyin{mean}: indipendente dall'area e meno fragile di \pyin{max}, che dipende da un
solo pixel e premia i box grandi per pura probabilità di contenere un picco isolato.

\subsection{\texttt{rank\_by\_saliency}}

\begin{idebox}
\begin{pycode}
scores = np.array([score_box(saliency_map, b, strategy) for b in boxes])
order = np.argsort(scores)[::-1]

return [
    {"index": int(i), "box": tuple(boxes[i]),
     "score": float(scores[i]), "rank": position + 1}
    for position, i in enumerate(order)
]
\end{pycode}
\end{idebox}

\begin{attenzione}
\textbf{Il campo \texttt{index} è essenziale.} La lista è ordinata per punteggio, quindi la
posizione nella lista \emph{non} corrisponde più al box originale. Senza \pyin{index} sarebbe
impossibile allineare i due ranking per il confronto: si finirebbe per confrontare oggetti
diversi credendo di confrontare gli stessi.
\end{attenzione}

\subsection{\texttt{rank\_by\_clicks}}

\begin{idebox}
\begin{pycode}
for click in clicks:
    containing = [
        i for i, (x1, y1, x2, y2) in enumerate(boxes)
        if x1 <= click["x"] < x2 and y1 <= click["y"] < y2
    ]
    if not containing:
        continue

    smallest = min(containing,
                   key=lambda i: (boxes[i][2] - boxes[i][0]) * (boxes[i][3] - boxes[i][1]))
    scores[smallest] += 1.0 / click.get("order", 1)
\end{pycode}
\end{idebox}

Due decisioni:

\begin{description}[leftmargin=1.2em,style=nextline]
\item[Click nello sfondo: scartati] Un click che non cade in nessun box indica qualcosa che il
detector non ha rilevato. Non c'è un oggetto a cui attribuirlo.

\item[Box sovrapposti: vince il più piccolo] I box si sovrappongono regolarmente (una persona e
il suo volto, un tavolo e gli oggetti sopra). Assegnare al box più piccolo assume che l'annotatore
intendesse l'elemento più specifico.
\end{description}

\begin{attenzione}
La seconda è un'\textbf{euristica}, non una certezza: in alcuni casi l'intenzione poteva essere
l'oggetto contenitore. Va dichiarata fra i limiti del metodo nel README.
\end{attenzione}

Il punteggio $1/\text{order}$ è lo stesso usato nella ground truth, per coerenza: primo click
1.0, secondo 0.5, terzo 0.33. La caduta è ripida di proposito --- la prima cosa notata è
qualitativamente diversa dalla quinta.

\subsection{\texttt{aligned\_scores}}

\begin{idebox}
\begin{pycode}
predicted = np.zeros(n_boxes, dtype=np.float64)
reference = np.zeros(n_boxes, dtype=np.float64)

for entry in predicted_ranking:
    predicted[entry["index"]] = entry["score"]
for entry in reference_ranking:
    reference[entry["index"]] = entry["score"]

return predicted, reference
\end{pycode}
\end{idebox}

\begin{attenzione}
\textbf{Il bug che questa funzione previene.} Passando direttamente i punteggi ordinati alle
metriche, si confronterebbe ``il primo classificato secondo il modello'' con ``il primo secondo
la ground truth'' --- ottenendo correlazione perfetta \emph{sempre}, indipendentemente da quali
oggetti siano.

\medskip
Un risultato troppo bello per essere vero, e senza alcun errore visibile.
\end{attenzione}

\section{\texttt{src/registry.py}}

\begin{idebox}
\begin{pycode}
_REGISTRY = {
    "Center Prior": (center_prior_saliency, False),
    "Spectral Residual": (spectral_residual_saliency, True),
    "Itti-Koch": (itti_koch_saliency, True),
}

AVAILABLE_MODELS = list(_REGISTRY.keys())


def compute_saliency(model_name, image_bgr):
    if model_name not in _REGISTRY:
        raise KeyError(
            f"Modello sconosciuto: '{model_name}'. "
            f"Disponibili: {', '.join(AVAILABLE_MODELS)}"
        )

    model_fn, needs_image = _REGISTRY[model_name]
    if needs_image:
        return model_fn(image_bgr)
    return model_fn(image_bgr.shape[:2])
\end{pycode}
\end{idebox}

Un dizionario che mappa il nome leggibile a una coppia \pyin{(funzione, vuole_immagine_intera)}.
Il flag booleano cattura l'unica differenza fra le firme: il center prior vuole la \emph{forma},
gli altri l'\emph{immagine}.

Il prefisso \pyin{_} in \pyin{_REGISTRY} segnala che è un dettaglio interno: chi usa il modulo
deve passare da \pyin{AVAILABLE_MODELS} e \pyin{compute_saliency}.

Il messaggio d'errore elenca i modelli disponibili invece di limitarsi a segnalare il fallimento.

\begin{nota}
\textbf{Il beneficio concreto.} Aggiungere un modello richiede di modificare \emph{soltanto}
questo file: una riga nel dizionario. L'app, i test e lo script di valutazione lo vedranno
comparire automaticamente, perché iterano su \pyin{AVAILABLE_MODELS}.
\end{nota}

\section{\texttt{src/visualization.py}}

\subsection{\texttt{saliency\_to\_heatmap}}

\begin{idebox}
\begin{pycode}
saliency_uint8 = (np.clip(saliency_map, 0.0, 1.0) * 255).astype(np.uint8)
heatmap_bgr = cv2.applyColorMap(saliency_uint8, cv2.COLORMAP_JET)
return cv2.cvtColor(heatmap_bgr, cv2.COLOR_BGR2RGB)
\end{pycode}
\end{idebox}

\pyin{cv2.applyColorMap} lavora solo su \pyin{uint8}, quindi la mappa $[0,1]$ va riscalata a
$[0,255]$. Il \pyin{np.clip} è una difesa: se un modello restituisse valori leggermente fuori
range, il cast produrrebbe wrap-around (255.5 diventa 0) con artefatti vistosi.

\pyin{COLORMAP_JET} è lo standard de facto nella letteratura sulla saliency: blu = basso, rosso =
alto. La conversione finale a RGB serve perché OpenCV produce BGR mentre Gradio si aspetta RGB.

\subsection{\texttt{overlay\_saliency}}

\begin{idebox}
\begin{pycode}
return cv2.addWeighted(image_rgb, 1.0 - alpha, heatmap_rgb, alpha, 0.0)
\end{pycode}
\end{idebox}

\pyin{cv2.addWeighted(src1, a, src2, b, gamma)} calcola \pyin{src1*a + src2*b + gamma}. Con
\pyin{a = 1-alpha} e \pyin{b = alpha} si ottiene una miscelazione lineare: \pyin{alpha=0} mostra
solo l'immagine, \pyin{alpha=1} solo la heatmap.

\subsection{\texttt{mark\_peak}}

\begin{idebox}
\begin{pycode}
annotated = image_rgb.copy()   # cv2.circle disegna in-place
py, px = np.unravel_index(np.argmax(saliency_map), saliency_map.shape)

cv2.circle(annotated, (int(px), int(py)), radius + 2, (0, 0, 0), thickness + 2)
cv2.circle(annotated, (int(px), int(py)), radius, (255, 255, 255), thickness)
\end{pycode}
\end{idebox}

Il \pyin{.copy()} è essenziale: senza, \pyin{cv2.circle} modificherebbe l'immagine del chiamante.

Vengono disegnati \textbf{due} cerchi concentrici, uno nero più spesso sotto e uno bianco più
sottile sopra: così il marcatore resta visibile sia su sfondi chiari sia su sfondi scuri.

\begin{attenzione}
\pyin{np.argmax} restituisce l'indice nell'array \emph{appiattito};
\pyin{np.unravel_index} lo riconverte in \pyin{(riga, colonna)}, cioè \pyin{(y, x)}. Ma
\pyin{cv2.circle} vuole il centro come \pyin{(x, y)}: l'ordine va invertito esplicitamente.
\end{attenzione}

\section{\texttt{app.py}}

\subsection{Principio architetturale}

Il file contiene \textbf{solo} l'interfaccia. Il calcolo vive in \pyin{src/registry.py}, la
visualizzazione in \pyin{src/visualization.py}.

\subsection{\texttt{prepare\_image}}

\begin{idebox}
\begin{pycode}
MAX_SIDE = 800

def prepare_image(image_rgb):
    height, width = image_rgb.shape[:2]
    scale = MAX_SIDE / max(height, width)

    if scale < 1.0:
        image_rgb = cv2.resize(image_rgb,
                               (int(width * scale), int(height * scale)),
                               interpolation=cv2.INTER_AREA)

    return image_rgb, cv2.cvtColor(image_rgb, cv2.COLOR_RGB2BGR)
\end{pycode}
\end{idebox}

Due responsabilità:

\begin{description}[leftmargin=1.2em,style=nextline]
\item[Ridimensionamento condizionale] \pyin{scale} si calcola sul lato \emph{maggiore}, così il
rapporto d'aspetto è preservato. La condizione \pyin{scale < 1.0} fa sì che le immagini piccole
non vengano \emph{ingrandite} --- interpolare non aggiunge informazione.

\item[Conversione RGB $\to$ BGR] Gradio consegna immagini in RGB, i modelli lavorano in BGR. La
conversione avviene in un punto unico, all'ingresso.
\end{description}

Il ridimensionamento è necessario per motivi pratici: Itti-Koch costruisce sette livelli di
piramide su tre canali più quattro orientazioni. Su una foto da smartphone il calcolo diventa
lento senza alcun beneficio, dato che la saliency è un fenomeno a bassa frequenza spaziale.

\subsection{\texttt{process}}

\begin{idebox}
\begin{pycode}
def process(image_rgb, model_name, alpha, show_peak):
    if image_rgb is None:
        return None, None, "Carica un'immagine per iniziare."

    image_rgb, image_bgr = prepare_image(image_rgb)
    saliency = compute_saliency(model_name, image_bgr)

    overlay = overlay_saliency(image_rgb, saliency, alpha=alpha)
    heatmap = saliency_to_heatmap(saliency)
    ...
\end{pycode}
\end{idebox}

Il controllo iniziale gestisce lo stato transitorio dell'interfaccia: all'apertura non c'è ancora
un'immagine. Senza di esso l'app solleverebbe un'eccezione visibile nel browser.

\subsection{Il selettore a radio button}

\begin{idebox}
\begin{pycode}
model_selector = gr.Radio(
    choices=AVAILABLE_MODELS,
    value=AVAILABLE_MODELS[0],
    label="Modello",
)
\end{pycode}
\end{idebox}

\pyin{gr.Radio} permette di selezionare \textbf{un solo} modello alla volta --- a differenza di
\pyin{gr.CheckboxGroup}, che ne permetterebbe più d'uno. Le scelte vengono direttamente da
\pyin{AVAILABLE_MODELS}: aggiungendo un modello al registro, compare qui senza toccare l'app.

\subsection{Collegamento degli eventi}

\begin{idebox}
\begin{pycode}
inputs = [image_input, model_selector, alpha_slider, peak_checkbox]
outputs = [overlay_output, heatmap_output, info_output]

run_button.click(process, inputs=inputs, outputs=outputs)
image_input.change(process, inputs=inputs, outputs=outputs)
model_selector.change(process, inputs=inputs, outputs=outputs)
alpha_slider.change(process, inputs=inputs, outputs=outputs)
peak_checkbox.change(process, inputs=inputs, outputs=outputs)
\end{pycode}
\end{idebox}

Ogni collegamento dichiara: \emph{quando accade questo evento, chiama \pyin{process} passando i
valori di \pyin{inputs}, e usa il risultato per aggiornare \pyin{outputs}}.

\begin{attenzione}
L'ordine di \pyin{inputs} deve corrispondere \textbf{esattamente} all'ordine dei parametri di
\pyin{process}. È posizionale, non per nome: invertire due voci non produce un errore, solo
comportamenti sbagliati.
\end{attenzione}

Oltre al pulsante esplicito, ogni controllo ricalcola automaticamente al cambiamento: cambiando
il radio button il risultato si aggiorna subito.

\section{\texttt{annotate.py}}

\subsection{\texttt{AnnotationSession}}

Classe che incapsula lo stato della sessione: quale immagine, quali click, quale annotatore.
Separata dall'interfaccia per lo stesso motivo dell'app --- così è testabile senza browser.

\subsubsection{\texttt{start} --- ripresa automatica}

\begin{idebox}
\begin{pycode}
self.store = AnnotationStore(annotator, self.annotations_dir)
done = self.store.annotated_images()

self.index = 0
while self.index < len(self.image_names) and self.image_names[self.index] in done:
    self.index += 1
\end{pycode}
\end{idebox}

Il ciclo salta le immagini già annotate \emph{da quell'annotatore}. Chi riapre il browser il
giorno dopo riprende esattamente da dove aveva lasciato, senza rifare il lavoro.

\subsubsection{\texttt{save\_and\_advance}}

\begin{idebox}
\begin{pycode}
if len(self.clicks) < MIN_CLICKS:
    return False, f"Servono almeno {MIN_CLICKS} click (ne hai {len(self.clicks)})."

height, width = self.current_image.shape[:2]
self.store.save_annotation(self.current_name, (height, width), self.clicks)

self.index += 1
self.clicks = []
self._load_current()
return True, "Salvato."
\end{pycode}
\end{idebox}

Restituisce una coppia \pyin{(successo, messaggio)} invece di sollevare eccezioni: in
un'interfaccia, un vincolo non rispettato è una condizione normale da comunicare all'utente, non
un errore di programmazione.

Notare l'ordine: si salva \emph{prima} di avanzare. Se il salvataggio fallisse, l'indice non
avanzerebbe e il lavoro non andrebbe perso.

\subsection{\texttt{draw\_clicks}}

\begin{idebox}
\begin{pycode}
for index, (x, y) in enumerate(clicks, start=1):
    cv2.circle(annotated, (x, y), 16, (0, 0, 0), -1)
    cv2.circle(annotated, (x, y), 14, (255, 80, 40), -1)

    text = str(index)
    (tw, th), _ = cv2.getTextSize(text, cv2.FONT_HERSHEY_SIMPLEX, 0.6, 2)
    cv2.putText(annotated, text, (x - tw // 2, y + th // 2), ...)
\end{pycode}
\end{idebox}

Il \textbf{numero} sopra ogni marcatore rende visibile l'\emph{ordine}, che è metà del dato
raccolto: senza indicazione visiva l'annotatore perde il conto di cosa ha già cliccato.

\pyin{cv2.getTextSize} misura il testo prima di disegnarlo, così da centrarlo esattamente sul
marcatore: \pyin{x - tw // 2} sposta a sinistra di metà larghezza.

\subsection{Il gestore dei click}

\begin{idebox}
\begin{pycode}
def on_click(event: gr.SelectData):
    if session.current_image is None:
        return None, session.status()
    x, y = event.index[0], event.index[1]
    ...
\end{pycode}
\end{idebox}

\pyin{gr.SelectData} è l'oggetto evento di Gradio per il click su un'immagine.
\pyin{event.index} contiene \pyin{[x, y]} --- colonna, riga: l'ordine opposto a quello di numpy.

\section{\texttt{scripts/run\_evaluation.py}}

Entry point che collega tutti i componenti.

\subsection{Il controllo di coerenza dimensionale}

\begin{idebox}
\begin{pycode}
if (height, width) != tuple(record["image_size"]):
    print(f"  ATTENZIONE: {image_name} e' {(height, width)} ma le annotazioni "
          f"sono per {tuple(record['image_size'])}. Saltata.")
    continue
\end{pycode}
\end{idebox}

\begin{attenzione}
\textbf{Perché è essenziale.} Le coordinate dei click sono relative alla risoluzione mostrata
all'annotatore. Se l'immagine su disco venisse ridimensionata dopo l'annotazione, i click
punterebbero a posizioni diverse --- e la valutazione produrrebbe numeri plausibili ma sbagliati,
senza alcun errore visibile. È il motivo per cui \pyin{image_size} viene salvato con i click.
\end{attenzione}

\subsection{Gestione dei NaN nelle medie}

\begin{idebox}
\begin{pycode}
values = np.array([r[metric] for r in model_rows], dtype=np.float64)
valid = values[~np.isnan(values)]

if valid.size == 0:
    entry[f"{metric}_mean"] = float("nan")
else:
    entry[f"{metric}_mean"] = round(float(valid.mean()), 4)
\end{pycode}
\end{idebox}

I \pyin{nan} restituiti dalle metriche non definite vengono esclusi: senza il filtro, un singolo
\pyin{nan} renderebbe \pyin{nan} l'intera colonna.

\subsection{Le due tabelle prodotte}

\begin{center}
\begin{tabular}{>{\ttfamily}l p{8.2cm}}
\toprule
\rowcolor{tblHead}\textcolor{white}{\textbf{File}} & \textcolor{white}{\textbf{Contenuto}} \\
\midrule
metrics\_per\_image.csv & Una riga per coppia (immagine, modello). Permette di individuare i casi
in cui un modello fallisce. \\
\rowcolor{tblAlt} metrics\_summary.csv & Media e deviazione standard per modello. È la tabella del
report finale. \\
\bottomrule
\end{tabular}
\end{center}

La deviazione standard non è decorativa: un modello con media alta ma varianza altissima è meno
affidabile di uno leggermente peggiore ma costante.

\section{\texttt{tests/test\_all.py}}

56 controlli. La strategia: ogni verifica usa casi in cui la risposta corretta è nota
\emph{per costruzione} o \emph{per ragionamento matematico}.

\begin{nota}
\textbf{Perché non basta confrontarsi con un'altra libreria.} Un test che verifica solo l'accordo
con \texttt{scipy} non dice se si è capita la metrica: se si sbaglia una convenzione (per esempio
l'ordine degli argomenti di KL) il confronto potrebbe comunque passare. I casi con verità nota
verificano il \emph{significato}, non solo la coincidenza numerica.
\end{nota}

\subsection{Esempi di verità note}

\begin{center}
\begin{tabular}{l l}
\toprule
\rowcolor{tblHead}\textcolor{white}{\textbf{Caso}} & \textcolor{white}{\textbf{Valore atteso}} \\
\midrule
Mappa confrontata con se stessa & CC $=1$, SIM $=1$, KL $=0$ \\
\rowcolor{tblAlt} Mappa e il suo opposto & CC $=-1$ \\
Center prior 240$\times$320 & picco esattamente in $(160, 120)$ \\
\rowcolor{tblAlt} Ordinamento invertito & $\rho = -1$ \\
Una coppia invertita su sei & $\tau = 2/3$ \\
\rowcolor{tblAlt} Una fissazione sul minimo & AUC $= 0.5$ (limite inferiore) \\
\bottomrule
\end{tabular}
\end{center}

\subsection{Il test che documenta un limite}

\begin{idebox}
\begin{pycode}
gray = cv2.cvtColor(img, cv2.COLOR_BGR2GRAY)
square = gray[y0:y1, x0:x1].mean()
mask = np.ones(gray.shape, bool)
mask[y0:y1, x0:x1] = False

check(abs(square - gray[mask].mean()) < 10,
      f"quadrato ({square:.0f}) e sfondo ({gray[mask].mean():.0f}) "
      f"isoluminanti: SR non puo' vederlo, Itti-Koch si'")
\end{pycode}
\end{idebox}

L'immagine di test è un quadrato \textbf{rosso} su sfondo \textbf{verde}. In scala di grigi il
quadrato ha luminanza media \textbf{94}, lo sfondo \textbf{98}: sono praticamente
\emph{isoluminanti}.

\begin{nota}
\textbf{Un risultato che vale come argomento all'orale.} Spectral Residual converte in scala di
grigi come primo passo, quindi per lui quel quadrato \emph{non esiste}: si distingue solo per
tinta. Itti-Koch, avendo canali di opponenza cromatica espliciti, lo trova senza difficoltà.

\medskip
Non è un bug ma una \textbf{limitazione strutturale}, e mostra che i modelli non sono ordinabili
su una scala unica: falliscono in modi diversi, su input diversi. Genera anche una previsione
verificabile sul dataset reale --- Spectral Residual dovrebbe andare peggio su scene dove
l'oggetto interessante si distingue per colore ma non per luminosità.
\end{nota}

\subsection{Isolamento tramite cartelle temporanee}

\begin{idebox}
\begin{pycode}
tmp = tempfile.mkdtemp()
try:
    ...
finally:
    shutil.rmtree(tmp)
\end{pycode}
\end{idebox}

Ogni test che tocca il disco lavora in una cartella temporanea, ripulita nel \pyin{finally}
(quindi anche se il test fallisce). I test non toccano mai i dati reali del progetto.

\section{Riepilogo}

\begin{center}
\begin{tabular}{l c l}
\toprule
\rowcolor{tblHead}\textcolor{white}{\textbf{Modello}} & \textcolor{white}{\textbf{Usa il contenuto?}} & \textcolor{white}{\textbf{Punto debole}} \\
\midrule
Center prior & no & ignora tutto \\
\rowcolor{tblAlt} Spectral Residual & sì & cieco al colore puro \\
Itti-Koch & sì & nessuna semantica \\
\bottomrule
\end{tabular}
\end{center}

Il limite comune a tutti e tre: \textbf{nessuno capisce cosa sia un oggetto}. Rispondono a
contrasto e statistica di basso livello. Un volto umano e una macchia colorata della stessa
dimensione e contrasto ricevono lo stesso trattamento --- mentre per un osservatore umano il
volto vince sempre.

È esattamente il divario che i modelli deep learning, addestrati su fissazioni umane reali,
dovrebbero colmare: l'estensione naturale del progetto.

\subsection{Domande di autoverifica}

\begin{nota}
Provare a rispondere a mente, senza rileggere:

\medskip
\begin{enumerate}
  \item Perché \pyin{center_prior_saliency} accetta \pyin{image_shape} e non l'immagine?
  \item Cosa succederebbe se Spectral Residual scartasse la fase e ne usasse una casuale?
  \item Perché in Itti-Koch si divide per l'intensità nel calcolo dei canali di colore?
  \item L'operatore $N(\cdot)$ amplifica o sopprime una mappa con molti picchi simili? Perché?
  \item Perché la ground truth ha due rappresentazioni invece di una?
  \item Come si giustifica il valore $\sigma = 5\%$ della larghezza?
  \item Perché la strategia \pyin{sum} favorisce sistematicamente gli oggetti grandi?
  \item Cosa succede se si passano alle metriche i due ranking già ordinati per punteggio?
  \item Con una sola fissazione, qual è il valore minimo possibile dell'AUC-Judd? Perché?
  \item Perché aggregare i click di tutti gli annotatori è diverso dal mediare le loro heatmap?
\end{enumerate}
\end{nota}
